\section{Quality and Validity of Generated Tests}

This section applies test validity --- whether a test suite actually verifies the intended behavior, not just whether it runs and passes --- to Part 1 and then to Part 2.

Part 1's backend testing was not limited to a single feature; pytest suites were generated after implementation across the application's routes, services, and models more broadly. The five evaluation dimensions used in this thesis --- coverage, correctness, completeness, an overgeneration check, and debugging capability --- were applied informally to these suites as well. The collections test-failure case (Section 4.8) sits outside these five dimensions entirely: the test itself was correct, but the fix that followed patched around the failure instead of resolving it, which is the distinct problem test validity was introduced to address.

Part 2's three-tier structure requires a pytest integration test for every new endpoint, a smoke test for every new component, and browser-based verification for anything touching the interface, so coverage no longer depended on whether an agent happened to generate tests for a given part of a feature. The requirement-derived testing rule targets validity rather than coverage: a test suite is no longer complete because it runs and passes, but only once every requirement bullet in the feature's plan has a corresponding test, covering a normal case, an error case, and any stated boundary condition --- a test checking something true but irrelevant to the requirement no longer counts. The permission-denial verification captured for the workspace roles feature --- a viewer blocked from creating a collection, a non-owner blocked from inviting a member --- is a validity concern specific to access-control logic that nothing in Part 1's feature set ever required testing for. And the reasoning behind the collections case, informal and applied once in Part 1, became a mandatory step in Part 2, required for every failing test regardless of how minor it looks.

Across both parts, test validity moved from something assessed informally, feature by feature, to a structural requirement checked the same way for every feature --- and from a suite that could pass while checking too little or checking the wrong thing, to one where a passing result is tied explicitly back to a stated requirement rather than to whatever the implementation happens to do.

